% Based on templates/tutorial.tex; omit dialogue and learner-check records.
% Source: LEC06 Image Recognition - Tutorial Problems.pdf, PDF p. 1, Q1--Q3.
% Key: LEC06 Image Recognition - Tutorial Solutions.pdf, PDF pp. 2--12.
% Session: 2026-10-27, teaching Week 11. All three questions answered and checked.

\exercisegroup{SC4061-TUT06}{Tutorial 06：Image Recognition}
\label{tutorial:SC4061-TUT06}

\noteintro{
{\small 2026-10-27（第十一教学周）。来源：\emph{LEC06 Image Recognition -- Tutorial Problems}，PDF p.~1，Questions 1--3；参考解答 \emph{Tutorial Solutions} pp.~2--12。题目、计数与 ROC 已核对。}
}

\exercisequestion{SC4061-TUT06-Q01}{Visual Words 与整图表示}

\textbf{对应讲义：}\relatedtopic{SC4061-L09-T03}{SIFT 描述子}；
\relatedtopic{SC4061-L09-T04}{BoVW 共享词典与直方图}。

\textbf{题目。}解释 Bag-of-Words（词袋模型）中 visual words（视觉词）的基本概念。

\begin{checkedsolution}
从训练图像提取 SIFT 局部描述子，在描述子空间中进行聚类。每个 cluster centroid（聚类中心）充当一个视觉词，代表一类常见的局部视觉模式，例如相似的边缘、角点或纹理。视觉词无需具有明确的物体语义。全部中心构成 visual vocabulary（视觉词表）：
\[
  \mathcal V=\{\bm v_1,\ldots,\bm v_K\},\qquad
  \bm v_j\in\mathbb R^{128}.
\]
这里的 $128$ 来自 SIFT 描述子的维数，$K$ 是聚类中心数，即词表大小。

表示一张图片时，将其每个局部描述子分配给最近的视觉词，再按词表顺序累计次数，得到 $K$ 维直方图。因此，\textbf{一个视觉词是 $128$ 维中心，整张图片的 BoW 表示是 $K$ 维向量}；图中关键点数量可以变化。
\end{checkedsolution}

\exercisequestion{SC4061-TUT06-Q02}{SPM 如何保留空间布局}

\textbf{对应讲义：}\relatedtopic{SC4061-L09-T04}{BoVW 的空间信息局限}。

\textbf{题目。}说明 Spatial Pyramid Matching（SPM，空间金字塔匹配）相比普通 Bag-of-Words 能改善场景分类表现的主要原因。

\begin{checkedsolution}
普通 BoW 在整张图上累计视觉词频数，丢失局部特征之间的空间布局。若两张图的视觉词种类及数量相同，即使各词的位置分布不同，全局直方图仍可相同。

SPM 把图像逐层划分为更细的子区域，\textbf{每个区域分别统计同一词表上的直方图}，并按固定区域顺序组合：
\[
  1\times1\quad\longrightarrow\quad2\times2\quad\longrightarrow\quad4\times4.
\]
粗层描述整体分布，细层描述模式所在的大致位置。三个层级分别有 $1$、$4$、$16$ 个区域；若每格直方图为 $K$ 维，直接拼接得到 $(1+4+16)K=21K$ 维表示。不同层级还可加权组合。

\begin{center}
  % Original schematic based on Tutorial 6 Solutions, PDF p. 8.
% Identical visual words and positions are reused at all three levels.
\begin{tikzpicture}[x=1cm,y=1cm,font=\small]
  \foreach \shift/\level/\n/\dimension in {0/0/1/1,4.2/1/2/4,8.4/2/4/16}{
    \begin{scope}[shift={(\shift,0)}]
      \fill[black!2] (0,0) rectangle (2.4,2.4);
      \ifnum\n>1
        \pgfmathtruncatemacro{\lastdivision}{\n-1}
        \foreach \k in {1,...,\lastdivision}{
          \pgfmathsetmacro{\division}{2.4*\k/\n}
          \draw[black!35,line width=0.4pt] (\division,0) -- (\division,2.4);
          \draw[black!35,line width=0.4pt] (0,\division) -- (2.4,\division);
        }
      \fi
      \draw[noteBlue,line width=0.7pt] (0,0) rectangle (2.4,2.4);
      \foreach \x/\y in {0.28/2.12,0.93/1.45,1.49/0.29,2.12/0.91}{
        \fill[noteBlue] (\x,\y) circle[radius=2pt];
      }
      \foreach \x/\y in {0.32/0.87,0.91/0.27,1.51/2.09,2.09/1.49}{
        \fill[ntuRed] (\x-0.055,\y-0.055) rectangle (\x+0.055,\y+0.055);
      }
      \foreach \x/\y in {0.29/1.46,0.93/2.12,1.48/0.92,2.10/0.28}{
        \fill[black!60] (\x,\y+0.075) -- (\x-0.065,\y-0.045)
          -- (\x+0.065,\y-0.045) -- cycle;
      }
      \node[anchor=south] at (1.2,2.52) {Level \level: $\n\times\n$};
      \node[anchor=north] at (1.2,-0.12) {拼接维数：$\dimension K$};
    \end{scope}
  }
  \node[anchor=north,font=\footnotesize,black!75] at (5.4,-0.59)
    {每格统计一个 $K$ 维直方图；相同形状与颜色表示相同视觉词。};
\end{tikzpicture}

  \par\smallskip
  \refstepcounter{figure}\label{fig:sc4061-tut06-spm}
  {\small 图~\thefigure：按参考解答 p.~8 的分层思想原创绘制。相同形状表示同类视觉词；每层观察同一组特征位置。}
\end{center}

因此，词频相同但\emph{区域分布不同}的场景可以得到不同表示，为分类增加布局线索。它保留的是分区粒度的空间信息，同一最细分区内部的排列仍可能无法区分；效果提升取决于这些布局线索对任务是否有用。
\end{checkedsolution}

\exercisequestion{SC4061-TUT06-Q03}{阈值评价与 ROC 曲线}

\textbf{对应讲义：}\relatedtopic{SC4061-L10-T02}{Precision、Recall 与 F1}。

\textbf{题目。}十个测试样本预测为正类的后验概率 $p_i$ 和真实类别如下。
(a) 在阈值 $0.8$ 下计算 precision、recall 和 F-measure；
(b) 绘制 ROC 曲线。

\begin{center}
\small
\setlength{\tabcolsep}{4pt}
\begin{tabular}{@{}l*{10}{c}@{}}
\toprule
样本 $i$ & 1 & 2 & 3 & 4 & 5 & 6 & 7 & 8 & 9 & 10 \\
\midrule
概率 $p_i$ & 0.95 & 0.93 & 0.87 & 0.852 & 0.851 & 0.85 & 0.76 & 0.53 & 0.43 & 0.25 \\
真实类别 & $+$ & $-$ & $-$ & $-$ & $+$ & $-$ & $-$ & $+$ & $-$ & $+$ \\
\bottomrule
\end{tabular}
\end{center}

\begin{checkedsolution}
\textbf{(a) 固定阈值。}约定 $p_i\ge t$ 时预测为正类。原表没有概率恰好等于 $0.8$ 的样本，因此此处用 $>$ 或 $\ge$ 得到相同结果。阈值 $t=0.8$ 时，样本 $1$--$6$ 判正，$7$--$10$ 判负：

\begin{center}
\begin{tabular}{@{}lcc@{}}
\toprule
真实类别 & 预测为正 & 预测为负 \\
\midrule
正类，共 4 个 & $TP=2$：$1,5$ & $FN=2$：$8,10$ \\
负类，共 6 个 & $FP=4$：$2,3,4,6$ & $TN=2$：$7,9$ \\
\bottomrule
\end{tabular}
\end{center}

Precision（精确率）衡量判正结果中真正例的比例；recall（召回率）衡量真实正类中被检出的比例：
\[
  P=\frac{TP}{TP+FP}=\frac{2}{6}=\boxed{\frac13},\qquad
  R=\frac{TP}{TP+FN}=\frac{2}{4}=\boxed{\frac12}.
\]
本题 F-measure 取两者的调和平均，即 $F_1$：
\[
  F_1=\frac{2PR}{P+R}
     =\frac{2TP}{2TP+FP+FN}
     =\frac{4}{4+4+2}=\boxed{\frac25}.
\]

\textbf{(b) 改变阈值，得到 ROC。}ROC 的横轴为 false positive rate（假正率），纵轴为 true positive rate（真正率）：
\[
  \boxed{FPR=\frac{FP}{FP+TN}=\frac{FP}{6}},\qquad
  \boxed{TPR=\frac{TP}{TP+FN}=\frac{TP}{4}=R}.
\]
两条坐标轴分别以\emph{全部真实负类}和\emph{全部真实正类}为分母。Precision 的分母是预测正类数，不能作为 ROC 横轴。

从 $t>0.95$ 开始，全部样本判负，曲线起点为 $(0,0)$。逐渐降低阈值，按表中分数降序把样本逐个纳入预测正类：遇到真实正类，$TP$ 加 $1$，\textbf{向上走 $1/4$}；遇到真实负类，$FP$ 加 $1$，\textbf{向右走 $1/6$}。当全部样本判正时，终点为 $(1,1)$。各步按降序连接，得到图~\ref{fig:sc4061-tut06-roc}。

\begin{center}
  % Original plot based on Tutorial 6 Solutions, PDF p. 12.
% Descending scores: .95, .93, .87, .852, .851, .85, .76, .53, .43, .25.
% Corresponding classes: +, -, -, -, +, -, -, +, -, + (P=4, N=6).
\begin{tikzpicture}[x=6.6cm,y=4.4cm,font=\small,>=Latex]
  \foreach \y in {0.25,0.5,0.75,1}{
    \draw[black!12,line width=0.4pt] (0,\y) -- (1,\y);
  }
  \foreach \x in {0.1666667,0.3333333,0.5,0.6666667,0.8333333,1}{
    \draw[black!12,line width=0.4pt] (\x,0) -- (\x,1);
  }
  \draw[->,black!70] (0,0) -- (1.08,0) node[right] {FPR};
  \draw[->,black!70] (0,0) -- (0,1.09) node[above] {TPR};
  \foreach \x/\tick in {0/0,0.1666667/{\frac16},0.3333333/{\frac13},
      0.5/{\frac12},0.6666667/{\frac23},0.8333333/{\frac56},1/1}{
    \draw[black!70] (\x,0) -- (\x,-0.016);
    \node[below=3pt,font=\footnotesize] at (\x,-0.016) {$\tick$};
  }
  \foreach \y/\tick in {0.25/{\frac14},0.5/{\frac12},0.75/{\frac34},1/1}{
    \draw[black!70] (0,\y) -- (-0.011,\y);
    \node[left=3pt,font=\footnotesize] at (-0.011,\y) {$\tick$};
  }
  \draw[noteBlue,line width=1.2pt]
    (0,0) -- (0,0.25) -- ({1/6},0.25) -- ({1/3},0.25)
    -- (0.5,0.25) -- (0.5,0.5) -- ({2/3},0.5)
    -- ({5/6},0.5) -- ({5/6},0.75) -- (1,0.75) -- (1,1);
  \foreach \x/\y in {0/0,0/0.25,0.1666667/0.25,0.3333333/0.25,
      0.5/0.25,0.5/0.5,0.6666667/0.5,0.8333333/0.5,
      0.8333333/0.75,1/0.75,1/1}{
    \fill[noteBlue] (\x,\y) circle[radius=1.7pt];
  }
  \draw[ntuRed,line width=0.6pt]
    ({2/3},0.54) -- ({2/3},0.63) -- (0.48,0.63) -- (0.48,0.66);
  \node[ntuRed,anchor=south,align=center,font=\footnotesize] at (0.48,0.67)
    {threshold $=0.8$\\$(\mathrm{FPR},\mathrm{TPR})=(\frac23,\frac12)$};
  \filldraw[draw=ntuRed,fill=white,line width=1pt]
    ({2/3},0.5) circle[radius=3.2pt];
  \fill[ntuRed] ({2/3},0.5) circle[radius=1.4pt];
\end{tikzpicture}

  \par\smallskip
  \refstepcounter{figure}\label{fig:sc4061-tut06-roc}
  {\small 图~\thefigure：由原题数据逐点计算并重绘的 ROC；突出点对应阈值 $0.8$。}
\end{center}

阈值 $0.8$ 的点为
\[
  (FPR,TPR)=\left(\frac{4}{6},\frac{2}{4}\right)
           =\boxed{\left(\frac23,\frac12\right)}.
\]
单个阈值决定一个点，扫描阈值得到整条曲线。计算时保留 $0.852>0.851>0.85$ 的原始精度，避免四舍五入制造并列分数、改变阶梯路径。参考解答 p.~12 用 \texttt{perfcurve(labels,scores,1)} 计算同样的 ROC 坐标。
\end{checkedsolution}

\textbf{一分钟回顾。}视觉词由局部描述子聚类得到，BoW 累计词频；SPM 分层、分区统计，补充粗略空间布局。固定阈值计算混淆矩阵与 $P,R,F_1$，改变阈值则得到以 $FPR$、$TPR$ 为坐标的 ROC。
